\markdownRendererDocumentBegin
\markdownRendererSectionBegin
\markdownRendererHeadingOne{Binary Lifting}\markdownRendererInterblockSeparator
{}Binary Lifting es una técnica de preprocesamiento para responder consultas de ancestros en árboles en tiempo logarítmico. Precalcula para cada nodo sus ancestros a distancias potencias de dos (\markdownRendererDollarSign{}2\markdownRendererCircumflex{}\markdownRendererLeftBrace{}0\markdownRendererRightBrace{}\markdownRendererDollarSign{}, \markdownRendererDollarSign{}2\markdownRendererCircumflex{}\markdownRendererLeftBrace{}1\markdownRendererRightBrace{}\markdownRendererDollarSign{}, \markdownRendererDollarSign{}2\markdownRendererCircumflex{}\markdownRendererLeftBrace{}2\markdownRendererRightBrace{}\markdownRendererDollarSign{}, \markdownRendererBackslash{}ldots) mediante programación dinámica, permitiendo alcanzar cualquier ancestro k-ésimo mediante descomposición binaria de k. Con complejidad \markdownRendererDollarSign{}O(n \markdownRendererBackslash{}log n)\markdownRendererDollarSign{} de preprocesamiento y \markdownRendererDollarSign{}O(\markdownRendererBackslash{}log n)\markdownRendererDollarSign{} por consulta, es fundamental para implementar Lowest Common Ancestor (LCA) eficientemente.\markdownRendererInterblockSeparator
{}\markdownRendererInputFencedCode{./_markdown_main/d19075fdc4e49aa779a39f8028919ef0.verbatim}{cpp}{cpp}
\markdownRendererSectionEnd \markdownRendererDocumentEnd